\subsection{Ambiente}
Para demostrar a comparação de performance foram gerados dados em ambiente controlado tanto temporal quanto em questão de \textit{hardware}, para garantir uma melhor veracidade dos dados, sendo realizado em uma máquina contendo:

\begin{itemize}
    \item Processador - AMD Ryzen 5 7535HS.
    \item Memória RAM - 8 \textit{gigabytes} DDR5 4800 MHz.
    \item Sistema operacional - Windows 10 PRO.
    \item Armazenamento - 512 GB SSD NVMe PCIe 4.0.
    \item PHP - 8.2.0
    \item SGBD - MySQL 8.0.0
\end{itemize}

Cada teste foi realizado em \textit{loops} de cem repetições ou de mil repetições, este último para processos mais rápidos. Cada teste foi adaptado para cada estrutura de cada um dos \textit{frameworks} mais difundidos no contexto deste trabalho. Para não haver influências terceiras de outras funcionalidades, foi utilizada uma série de abstrações feitas em PHP nativo sem quaisquer resquício de ferramentas dos \textit{frameworks} para realizar o teste, à exceção do \textit{script} sujeito à análise. Os testes foram realizados em uma base de dados contendo as seguintes quantidades de registros:

\begin{itemize}
    \item Tabela de usuários - 1000.
    \item Tabela de postagens - 10000.
    \item Tabela de empresas - 100.
    \item Tabela relacionando usuários e postagens - 100000.
\end{itemize}